\begin{solapartat}
\punts{0,25} $E = mc^2$

\punts{0,1} $E = 9,11\times 10^{-31}(3,00\times 10^8)^2 = 8,20\times 10^{-14}\ \mathrm{J}$

\punts{0,1} $E = 8,20\times 10^{-14}\ \mathrm{J}\ \dfrac{1\ \mathrm{eV}}{1,602\times 10^{-19}\ \mathrm{J}}\ \dfrac{1\ \mathrm{MeV}}{10^6\ \mathrm{eV}} = 0,512\ \mathrm{MeV}$

\punts{0,2} L'energia mínima d'un fotó per materialitzar-se en un parell electró-positró és el doble de l'energia en repòs del positró, ja que l'energia en repòs de l'electró és la mateixa que la del positró: $E = 2\times 0,512\ \mathrm{MeV} = 1,02\ \mathrm{MeV}$

\punts{0,2} La freqüència d'aquest fotó és:
\[ E_\text{fotó} = hf \Rightarrow f = \frac{2\times 8,20\times 10^{-14}}{6,63\times 10^{-34}} = 2,47\times 10^{20}\ \mathrm{Hz}. \]

\punts{0,2} La longitud d'ona d'aquest fotó és:
\[ \lambda = \frac{c}{f} = \frac{3,00\times 10^8}{2,47\times 10^{20}} = 1,21\times 10^{-12}\ \mathrm{m} \]

\punts{0,2} La quantitat de moviment d'aquest fotó és:
\[ p = \frac{h}{\lambda} = \frac{E_\text{fotó}}{c} = \frac{6,63\times 10^{-34}}{1,21\times 10^{-12}} = 5,47\times 10^{-22}\ \mathrm{kg\,m\,s^{-1}} \]
\end{solapartat}

\begin{solapartat}
\punts{0,25} En la desintegració $\beta^+$ es crea un positró:
\[ {}^{1}_{1}p^+ \to {}^{1}_{0}n + {}^{0}_{+1}e^+ + {}^{0}_{0}\nu \]

\punts{0,25} En la desintegració $\beta^-$ la partícula que es crea que és antimatèria és l'antineutrí.
\[ {}^{1}_{0}n \to {}^{1}_{1}p^+ + {}^{0}_{-1}e^- + {}^{0}_{0}\bar{\nu} \]

\punts{0,25} En el procés d'anihilació positró-electró, es crea un parell de fotons que viatgen en sentits oposats per a conservar la quantitat de moviment.

\punts{0,25} La càrrega inicial, abans d'anihilar-se, és zero, ja que abans de l'anihilació tenim $e^+ + e^-$, la càrrega positiva del positró té el mateix valor absolut que la negativa de l'electró, per tant, sumades ens dona una càrrega neta nul·la.

\punts{0,25} També és zero la final, després de l'anihilació: $e^+ + e^- \longrightarrow \gamma + \gamma$, tenim dos fotons i els fotons tenen càrrega nul·la. També sabem que en tot procés la carrega neta es conserva, per tant si inicialment era zero, al final també ho ha de ser.
\end{solapartat}
